\section{Mid-Term Examination Suite}

\ExamHeader{Mid-Term Examination Suite}{Set A (Model Paper 1)}{90 Minutes}{30 Marks}{-0.25 Marks}

\noindent
\textbf{General Instructions:}
\begin{enumerate}[topsep=2pt,itemsep=1pt]
    \item This paper contains \textbf{30 Multiple-Choice Questions (MCQs)} carrying \textbf{1 mark each} (Total = 30 Marks).
    \item \textbf{10 Questions from Unit 1} (Matrices), \textbf{10 Questions from Unit 2} (Differential Calculus), and \textbf{10 Questions from Unit 3} (Integral Calculus).
    \item Each incorrect response incurs a deduction of \textbf{0.25 marks}.
\end{enumerate}
\vspace{3mm}

\subsection*{Section A: Unit 1 --- Matrix Methods and Linear Systems (Questions 1 to 10)}
\mcqitem{1}{The rank of an identity matrix $I_n$ of order $n$ is:}{0}{1}{n}{$n-1$}
\mcqitem{2}{If $A$ is a $3 	imes 3$ matrix with $\det(A) = 4$, then $\det(2A)$ is:}{8}{16}{32}{64}
\mcqitem{3}{The sum of eigenvalues of $A = \begin{bmatrix} 4 & 1 \\ 2 & 3 \end{bmatrix}$ is:}{7}{10}{12}{5}
\mcqitem{4}{If $A^T = -A$, the diagonal elements of $A$ are all:}{1}{-1}{0}{Purely imaginary}
\mcqitem{5}{A system $Ax=b$ has a unique solution if and only if:}{$\rho(A) < \rho([A|b])$}{$\rho(A) = \rho([A|b]) = n$}{$\rho(A) < n$}{$\det(A) = 0$}
\mcqitem{6}{If $\lambda$ is an eigenvalue of $A$, then an eigenvalue of $A^2$ is:}{$\lambda$}{$\lambda^2$}{$2\lambda$}{$\sqrt{\lambda}$}
\mcqitem{7}{The product of eigenvalues of $A$ equals:}{$\operatorname{tr}(A)$}{$\det(A)$}{$\operatorname{rank}(A)$}{1}
\mcqitem{8}{If $A$ is orthogonal, then $A^{-1}$ equals:}{$A$}{$A^T$}{$-A$}{$I$}
\mcqitem{9}{A square matrix $A$ is non-singular if:}{$\det(A) = 0$}{$\det(A) \ne 0$}{$\operatorname{tr}(A) = 0$}{$\rho(A) = 0$}
\mcqitem{10}{The Cayley--Hamilton theorem states that every square matrix satisfies its own:}{Linear system}{Characteristic equation}{Adjugate matrix}{Inverse}

\subsection*{Section B: Unit 2 --- Differential Calculus and Applications (Questions 11 to 20)}
\mcqitem{11}{The $n$-th derivative of $e^{3x}$ is:}{$3 e^{3x}$}{$3^n e^{3x}$}{$n 3 e^{3x}$}{$e^{3x}/3^n$}
\mcqitem{12}{If $y = \sin(2x)$, then $y_n$ equals:}{$2^n \sin(2x + n\pi/2)$}{$2^n \cos(2x)$}{$2^n \sin(2x)$}{$\sin(2x + n\pi)$}
\mcqitem{13}{In Rolle's theorem for $f(x)$ on $[a,b]$, the derivative $f'(c) = 0$ holds for:}{$c \in [a,b]$}{$c \in (a,b)$}{$c = a$}{$c = b$}
\mcqitem{14}{The value of $\lim_{x\to 0} \frac{e^x - 1}{x}$ is:}{0}{1}{$\infty$}{$e$}
\mcqitem{15}{The Maclaurin series of $\cos x$ contains only:}{Odd powers of $x$}{Even powers of $x$}{Negative powers of $x$}{Fractional powers}
\mcqitem{16}{According to Lagrange's MVT, $f(b)-f(a) = $}{$(b-a)f''(c)$}{$(b-a)f'(c)$}{$f'(c)/(b-a)$}{0}
\mcqitem{17}{The $n$-th derivative of a product $u v$ is given by:}{Taylor's theorem}{Leibniz's theorem}{Euler's theorem}{Rolle's theorem}
\mcqitem{18}{The value of $\lim_{x\to 0} \frac{\tan x - x}{x^3}$ is:}{1/3}{-1/3}{1/6}{0}
\mcqitem{19}{If $f''(c) < 0$ at a critical point where $f'(c) = 0$, then $x=c$ is a:}{Local Minimum}{Local Maximum}{Inflection point}{Saddle point}
\mcqitem{20}{The $n$-th derivative of $\ln(x)$ is:}{$(-1)^{n-1}(n-1)!/x^n$}{$n!/x^n$}{$1/x^n$}{$(-1)^n n!/x^{n+1}$}

\subsection*{Section C: Unit 3 --- Fundamentals of Integral Calculus (Questions 21 to 30)}
\mcqitem{21}{The value of $\int_0^{\pi/2} \frac{\sin^3 x}{\sin^3 x + \cos^3 x}\,dx$ is:}{$\pi$}{$\pi/2$}{$\pi/4$}{0}
\mcqitem{22}{If $f(x)$ is an odd function, then $\int_{-a}^a f(x)\,dx$ equals:}{$2\int_0^a f(x)dx$}{0}{$a$}{$-a$}
\mcqitem{23}{The integral $\int \ln x\,dx$ equals:}{$x\ln x - x + C$}{$1/x + C$}{$x\ln x + x + C$}{$x^2/2 + C$}
\mcqitem{24}{The King's property states that $\int_a^b f(x)\,dx$ equals:}{$\int_a^b f(a-x)dx$}{$\int_a^b f(a+b-x)dx$}{$\int_a^b f(b-x)dx$}{$\int_a^b f(x/2)dx$}
\mcqitem{25}{The antiderivative of $\frac{1}{1+x^2}$ is:}{$\sin^{-1}x$}{$\tan^{-1}x$}{$\ln(1+x^2)$}{$\sec^{-1}x$}
\mcqitem{26}{The priority order for integration by parts is governed by:}{BODMAS}{ILATE}{L'Hopital}{Euler}
\mcqitem{27}{The value of $\int_0^1 x e^x\,dx$ is:}{1}{$e - 1$}{$e$}{0}
\mcqitem{28}{The integral $\int_{-1}^1 x^3 \sqrt{1-x^2}\,dx$ is:}{1}{0}{$\pi/2$}{2}
\mcqitem{29}{The integral of $\sec^2 x$ is:}{$\tan x + C$}{$\sec x + C$}{$\ln|\sec x| + C$}{$\tan^2 x + C$}
\mcqitem{30}{If $f(2a-x) = f(x)$, then $\int_0^{2a} f(x)\,dx$ equals:}{$2\int_0^a f(x)dx$}{0}{$\int_0^a f(x)dx$}{$4\int_0^a f(x)dx$}

\newpage
\subsection*{Complete Answer Key \& Technical Rationales}

\begin{table}[htbp]
\centering
\caption{\textbf{Answer Key \& Explanations (Questions 1 to 30)}}
\vspace{2mm}
\begin{tabularx}{\textwidth}{l c X l c X}
\toprule
\textbf{Q\#} & \textbf{Ans} & \textbf{Technical Rationale} & \textbf{Q\#} & \textbf{Ans} & \textbf{Technical Rationale} \\
\midrule
1 & \textbf{(c)} & All $n$ columns of $I_n$ are linearly independent pivots. & 16 & \textbf{(b)} & Definition of Lagrange's Mean Value Theorem. \\
2 & \textbf{(c)} & $\det(kA) = k^n\det(A) = 2^3(4) = 32$. & 17 & \textbf{(b)} & Leibniz theorem gives product expansion $(uv)_n$. \\
3 & \textbf{(a)} & Sum of eigenvalues equals trace $= 4 + 3 = 7$. & 18 & \textbf{(a)} & Series expansion: $\tan x = x + x^3/3 + \dots \implies 1/3$. \\
4 & \textbf{(c)} & $a_{ii} = -a_{ii} \implies 2a_{ii} = 0 \implies a_{ii} = 0$. & 19 & \textbf{(b)} & Second derivative test for local maximum. \\
5 & \textbf{(b)} & Rouch\'e--Capelli theorem condition for unique solution. & 20 & \textbf{(a)} & Standard $n$-th derivative formula for $\ln x$. \\
6 & \textbf{(b)} & Spectral property: $\lambda(A^k) = \lambda^k$. & 21 & \textbf{(c)} & King's property yields $2I = \pi/2 \implies I = \pi/4$. \\
7 & \textbf{(b)} & Product of all eigenvalues is identically equal to $\det(A)$. & 22 & \textbf{(b)} & Odd symmetry over $[-a, a]$ vanishes. \\
8 & \textbf{(b)} & $A A^T = I \implies A^{-1} = A^T$. & 23 & \textbf{(a)} & Integration by parts $\int 1 \cdot \ln x dx = x\ln x - x + C$. \\
9 & \textbf{(b)} & Non-singular matrices have non-zero determinant. & 24 & \textbf{(b)} & Fundamental reflection identity. \\
10 & \textbf{(b)} & $p(A) = O$ for characteristic polynomial. & 25 & \textbf{(b)} & Standard inverse trigonometric integral. \\
11 & \textbf{(b)} & Standard successive derivative: $(e^{ax})_n = a^n e^{ax}$. & 26 & \textbf{(b)} & ILATE: Inverse, Log, Algebraic, Trig, Exp. \\
12 & \textbf{(a)} & Standard formula for $n$-th derivative of $\sin(ax+b)$. & 27 & \textbf{(a)} & $[x e^x - e^x]_0^1 = (e - e) - (0 - 1) = 1$. \\
13 & \textbf{(b)} & Hypothesis guarantees $c$ in open interval $(a,b)$. & 28 & \textbf{(b)} & Integrand is odd, interval is symmetric $[-1, 1]$. \\
14 & \textbf{(b)} & L'H\^opital's rule: $\lim e^x / 1 = 1$. & 29 & \textbf{(a)} & Standard derivative: $d/dx(\tan x) = \sec^2 x$. \\
15 & \textbf{(b)} & $\cos(-x) = \cos x$ is even, so series has only even powers. & 30 & \textbf{(a)} & Standard definite integral reduction property. \\
\bottomrule
\end{tabularx}
\end{table}
